\subsection{以 a 为底的展开}

最重要的基序列是那些 \(d_{n}=a^{n}\) 的序列，其中 a 是整数 \textgreater1；这时 a 称为相应展开的底数（或简称底）。对于数值计算，使用以 10 为底的展开，称为十进展开；以 2 为底（二进展开）与以 3 为底（三进展开）的展开也常被使用。

为了表示不足近似 \(r_n\) （对数 \(x \geqslant 0\) ，在其以 \(a\) 为底的展开中），采用下述记号：每个整数 \(u\) （满足 \(0 \leqslant u \leqslant a - 1\) ）用一个特定的记号表示；如果

\[
r _ {n} = p _ {0} + \sum_ {k = 1} ^ {n} \frac {u _ {k}}{d _ {k}},
\]

我们首先借助这些记号写出以 \(a\) 为底的整数 \(p_{0}=[x]\geqslant0\) 的表示（集合论，第三章，§ 5，第 7 节），然后放一个点（``小数点''），并写出

相继地表示数 \(u_1, u_2, \ldots, u_n\) 的记号。如果 \(S\) 是这样得到的符号，习惯上（滥用语言）写 \(x = S \ldots\) 。应当一劳永逸地理解：这样的关系式只是指示右端是 \(x\) 的精确到 \(1/a^n\) 的不足近似的一种缩写方法。

对于负数，通行习惯不同：我们用上述记号写出 \(x' = -x > 0\) 的一个近似，并在其前放置符号``-''；因此这样记的是 \(x\) 的精确到 \(1/a^n\) 的过剩近似。

这个用法对数值计算有其不便之处。在负对数的记号中，采用与正数相同的记号，而在整数部分上放一横杠，以指示它等于所写数的相反数。

